\chapter{Fichiers, commandes et captures}

\section{Fichiers principaux}

\small
\begin{longtable}{p{0.43\textwidth}p{0.49\textwidth}}
\caption{Inventaire fonctionnel}\label{tab:fichiers}\\
\toprule
Fichier & Fonction \\
\midrule
\endfirsthead
\toprule
Fichier & Fonction \\
\midrule
\endhead
\fichier{arianes/pipelines/xmi_to_json.py} & Conversion XMI vers JSON. \\
\fichier{arianes/graph/uml_builder.py} & Construction du graphe et indexation documentaire. \\
\fichier{arianes/retrieval.py} & Routage et retrieval UML, métier et documentaire. \\
\fichier{arianes/agent.py} & Génération contrôlée multi-fournisseur. \\
\fichier{api.py} & API, authentification, conversations et messages. \\
\fichier{index.html} & Interface du chatbot. \\
\fichier{arianes/pipelines/extraction_metier.py} & Extraction contrôlée d'une fiche. \\
\fichier{arianes/pipelines/extraction_globale.py} & Extraction des 29 fiches avec reprise. \\
\fichier{arianes/pipelines/validation_metier.py} & Contrôle gradué V3. \\
\fichier{arianes/pipelines/recontrole_metier.py} & Recontrôle hors ligne des checkpoints. \\
\fichier{arianes/pipelines/generer_pilotes.py} & Génération d'un pilote par fiche. \\
\fichier{arianes/pipelines/selection_prudente.py} & Sélection documentaire globale. \\
\fichier{arianes/graph/metier_builder.py} & Import protégé dans Neo4j. \\
\fichier{arianes/evaluation.py} & Scénarios d'évaluation fonctionnelle. \\
\bottomrule
\end{longtable}
\normalsize

\section{Commandes utiles}

\begin{lstlisting}[language=bash,caption={Tests et lancement de l'application}]
python -m unittest discover -s tests -v
python -m uvicorn api:app --host 127.0.0.1 --port 8000
\end{lstlisting}

\begin{lstlisting}[language=bash,caption={Extraction et contrôle métier}]
python -m arianes.pipelines.extraction_globale
python -m arianes.pipelines.recontrole_metier --output-dir project_data/metier/corrections_v3
python -m arianes.pipelines.generer_pilotes
python -m arianes.pipelines.selection_prudente
python -m arianes.graph.metier_builder project_data/metier/pilote_global_documentaire.json
\end{lstlisting}

\section{Requêtes Neo4j}

\begin{lstlisting}[language={},caption={Visualisation du pilote métier}]
MATCH (c:ConceptMetier)
OPTIONAL MATCH (c)-[r:RELATION_METIER]->(n:ConceptMetier)
RETURN c, r, n;
\end{lstlisting}

\begin{lstlisting}[language={},caption={Visualisation de la classe Ligne}]
MATCH (c:Classe {nom: "Ligne"})
OPTIONAL MATCH (c)-[r]-(n)
RETURN c, r, n
LIMIT 40;
\end{lstlisting}

\section{Captures à ajouter}

Placer les images dans le dossier \code{figures/} avec les noms exacts indiqués dans \code{figures/README.md}. Le document compile sans elles et affiche des cadres réservés. Toute clé, tout mot de passe, cookie, jeton, identifiant personnel ou adresse interne doit être masqué avant insertion.
